\chapter{Introduction}
\label{chap:introduction}

\section{Context and motivation}

An autonomous vehicle does not make decisions from a static picture of the road. It repeatedly receives observations from perception, updates its estimate of the surrounding scene and selects an action while other road users continue to move. Motion forecasting is the component that bridges perception and planning: from an observed history, it estimates where nearby vehicles, cyclists and pedestrians may travel during a future horizon. A planner can then reason about potential conflicts before they occur. Forecasting is therefore not simply a curve-fitting exercise. It must represent interaction, road geometry, uncertainty and timing in a form that can be queried repeatedly under real-time constraints.

The future of a road user is not unique. A vehicle approaching a junction may turn, proceed or yield; a pedestrian may cross or wait. A useful predictor consequently returns several trajectories and an associated confidence for each one. This property is referred to as \emph{multimodality}: one observed history corresponds to multiple physically and socially plausible futures. Forecasting must also account for interactions. A target vehicle's future depends on relative position, heading, right-of-way and the likely actions of neighbouring agents. High-definition map elements constrain these possibilities further. The combined problem is difficult because a model must distinguish plausible alternatives without averaging them into an implausible path.

Most public benchmarks formulate motion forecasting as independent snapshots with a fixed observation duration. This formulation is valuable for standardised comparison, but it differs from a deployed system in which predictions are refreshed continuously. Treating each refresh as an unrelated scene duplicates computation and discards useful historical state. RealMotion addressed this mismatch by carrying context and prior trajectories across consecutive scenes \cite{song2024realmotion}. SEAM subsequently introduced endpoint-aware streaming \cite{prutsch2026seam}. SHARP further focused on agents with heterogeneous observation lengths: it processes non-overlapping short windows and transfers latent context between matched instances \cite{prutsch2026sharp}. These methods shift the research question from ``How accurately can a fixed history be encoded?'' to ``How can useful state be accumulated without weakening accuracy when observations arrive incrementally?''

SHARP is an appropriate basis for this dissertation because its central contribution is both practically relevant and sufficiently compact for controlled architectural investigation. The published model contains approximately 4.6 million parameters and reports competitive Argoverse 2 performance with low inference latency \cite{prutsch2026sharp}. More importantly, its streaming innovation is localised: short observation windows, instance-aware context transfer, target-centric context and trajectory relay are identifiable mechanisms that can be preserved while attention or temporal encoding is modified. This makes it possible to ask whether a proposed improvement helps forecasting accuracy rather than merely replacing the problem formulation.

\section{Problem statement}

Recent sequence models offer several plausible ways to improve SHARP. Transformer attention can represent direct pairwise interaction, but its logits may become excessively sharp when query and key magnitudes grow. Query--key normalisation (QKNorm) removes that magnitude dependence and learns the logit scale explicitly \cite{henry2020qknorm}. Relative-geometry bias can expose displacement, range and heading differences directly to scene attention rather than requiring them to be reconstructed from token content. Prediction uncertainty can also inform how widely a later streaming window searches around each previously predicted endpoint.

Mamba provides another alternative. It is a selective state-space model (SSM) whose recurrent parameters depend on the input, allowing information to be retained or forgotten while scaling linearly with sequence length \cite{gu2024mamba}. DeMo demonstrated that Mamba can be useful in motion forecasting when it is assigned to semantically ordered agent histories and future-state sequences \cite{zhang2024demo}. This does not imply that inserting Mamba anywhere in SHARP will improve it. A state-space scan assumes an order with meaningful recurrence. The ordering of pooled agent and lane tokens is not equivalent to time, and SHARP's ten-step local history is much shorter than the sequences for which Mamba is normally motivated. Placement and residual strength therefore require empirical study.

The practical problem addressed here is to determine which local architectural changes improve SHARP under controlled training, while preserving its defining short-window streaming behaviour. A second problem is methodological: long multi-GPU experiments are vulnerable to interruption, distributed-training faults and subtle configuration drift. Conclusions are only credible if each comparison records the actual global batch, schedule, checkpoint-selection rule and evaluation split, and if incomplete runs are not mistaken for final results.

\section{Research aim and questions}

The aim of this dissertation is to evaluate targeted attention, context and state-space modifications to SHARP for multimodal motion forecasting, using reproducible controlled ablations on Argoverse 2 while retaining the original streaming architecture.

The investigation is organised around four research questions:

\begin{enumerate}
    \item \textbf{RQ1:} How closely can the released SHARP model be reproduced on Argoverse 2, and which differences between the article, public configuration and local distributed environment affect that comparison?
    \item \textbf{RQ2:} Which of a broad set of local architecture modifications show evidence of improved displacement accuracy under an identical shortened training schedule?
    \item \textbf{RQ3:} Does changing the attention operator from standard multi-head attention to QKNorm, Talking-Heads attention, or their combination improve SHARP under a full 80-epoch controlled study?
    \item \textbf{RQ4:} Does Mamba benefit SHARP more as a residual chronological agent-history module than as a scene-token addition or a complete replacement for temporal attention?
\end{enumerate}

The final three-run study described in Chapters~\ref{chap:methodology} and \ref{chap:results} is a confirmatory extension of these questions. It tests whether individually promising changes remain beneficial when combined at the full article schedule. Its metric fields are reserved until validated output is available.

\section{Objectives}

The work was divided into the following objectives:

\begin{enumerate}
    \item establish an evidence-backed description of the SHARP article and released implementation, including data representation, streaming state, optimisation and evaluation;
    \item reproduce SHARP on Argoverse 2 and retain comparable validation metrics, logs, checkpoints and environment manifests;
    \item implement a controlled 20-epoch screen of eight proposed accuracy modifications and one complete temporal Mamba replacement, each isolated from the others;
    \item implement an 80-epoch attention study in which only the attention operator changes;
    \item compare two residual Mamba placements and one complete temporal-attention replacement without removing SHARP's instance-aware context stream;
    \item design a final full-length comparison combining only changes supported by the preceding evidence;
    \item make long-running experiments recoverable through per-epoch checkpoints, bounded preflight tests, warning scans and compact GitHub artifacts; and
    \item report negative as well as positive results, distinguishing a screening result, a best checkpoint, a final-epoch metric and an article value.
\end{enumerate}

\section{Contributions}

The principal contribution is not a claim that a single new layer universally improves motion prediction. It is a controlled analysis of where modest changes help or harm a compact streaming predictor. The dissertation makes five specific contributions.

First, it documents the article--code alignment required to train an article-schedule SHARP baseline. The released Hydra configuration and the article supplement do not use identical epoch and learning-rate values. The final protocol pins the released architecture but explicitly applies the article's 80-epoch schedule, rather than describing the repository defaults as the article setup.

Second, it provides a ten-variant screening study under identical data, seed, global batch and optimisation controls. The result narrows a broad design space: uncertainty-aware target context and relative-geometry bias improve minADE6 slightly, whereas several intuitively attractive additions do not.

Third, it provides a full-length controlled attention comparison. QKNorm improves the best minADE6 over the local MHA control, while Talking-Heads attention and the combined operator do not. This prevents the final design from being selected by architectural novelty alone.

Fourth, it isolates Mamba placement. Scene-token Mamba is applied after chronological detail has already been pooled; residual temporal-agent Mamba is applied within each ordered history; full replacement removes temporal attention entirely. Their outcomes support a restrained conclusion: the chronological residual placement is the strongest Mamba formulation tested, but the unmodified SHARP reproduction remains more accurate.

Fifth, it develops a reproducibility and recovery framework suitable for multi-day training. The framework records exact run directories, emits a last checkpoint every epoch, resumes incomplete variants, evaluates completed checkpoints separately, scans warnings and publishes small textual and graphical artifacts while retaining large checkpoints locally.

\section{Scope and exclusions}

The quantitative dissertation scope is SHARP on Argoverse 2, with secondary Lab 2 validation evidence on Argoverse 1 and nuScenes used only to discuss generalisation and reproduction gaps. Results from the earlier single-A4000 laboratory stage are excluded, as required by the project scope. SEAM is discussed as related work because endpoint-aware streaming helps locate SHARP within the literature; no SEAM experiment is reported. DeMo is used to explain why Mamba is relevant to motion forecasting, but the Mamba modifications developed here remain SHARP ablations rather than DeMo reproductions.

The dissertation does not claim closed-loop vehicle safety, planning improvement or deployment in CARLA/ROS2. These were considered in the initial project planning, but the completed technical evidence is a controlled forecasting study. Open-loop benchmark metrics are useful but cannot establish collision avoidance, calibration under distribution shift, or downstream planning value by themselves.

\section{Dissertation structure}

Chapter~\ref{chap:background} formulates multimodal forecasting, defines the metrics and reviews recurrent, vectorised, attention-based and streaming methods. Chapter~\ref{chap:methodology} details SHARP, the datasets, controlled training protocol, nine screening modifications, attention variants and Mamba placements. Chapter~\ref{chap:results} presents the completed Lab 2 and Lab 3 results, analyses cross-study consistency, and reserves the final confirmatory cells. Chapter~\ref{chap:conclusion} answers the research questions, states limitations and identifies the next experiments that can be justified by the evidence. The appendix records configurations, provenance and the result-insertion protocol.
